\section{PERFIL PROFISSIONAL DO EGRESSO}

Nos termos da \textbf{Resolução CNE/CES nº 2, de 24 de abril de 2019}, e da \textbf{Resolução CNE/CES nº 1, de 26 de março de 2021}, que instituem as Diretrizes Curriculares Nacionais para os cursos de Graduação em Engenharia, observa-se, especialmente, o disposto no Art. 3º, o qual estabelece que o perfil do egresso deve contemplar determinadas competências e características fundamentais à formação do engenheiro, a saber:

\begin{adjustwidth}{4cm}{0cm}
\it \footnotesize \setstretch{.5}

	O perfil do egresso do curso de graduação em Engenharia deve compreender, entre outras, as seguintes características:
	
	\begin{enumerate}[label=\Roman* -]
		\item visão holística e humanista, postura crítica, reflexiva, criativa, cooperativa e ética, sustentada por sólida formação técnico-científica;
		\item capacidade para pesquisar, desenvolver, adaptar e empregar novas tecnologias, atuando de modo inovador e empreendedor;
		\item habilidade para reconhecer necessidades dos usuários, formular, analisar e solucionar problemas de Engenharia de forma criativa;
		\item adoção de perspectivas multidisciplinares e transdisciplinares na prática profissional;
		\item consideração de aspectos globais, políticos, econômicos, sociais, ambientais, culturais, de segurança e de saúde no trabalho;
		\item atuação isenta e comprometida com a responsabilidade social e o desenvolvimento sustentável.
	\end{enumerate}
	
\end{adjustwidth}

Diante disso, o perfil esperado do futuro \NOMEPROFISSIONAL\ deve contemplar, obrigatoriamente, as seguintes características formativas:

\begin{itemize}
	\item Formação técnico-científica sólida, orientada por uma visão humanista, ética e crítica da realidade.
	\item Capacidade de pesquisar, desenvolver, adaptar e implementar tecnologias, com postura inovadora e empreendedora.
	\item Competência para identificar demandas, analisar contextos e formular soluções para problemas de Engenharia, empregando criatividade e rigor técnico.
	\item Integração de conhecimentos de distintas áreas, adotando abordagens multidisciplinares e transdisciplinares.
	\item Avaliação permanente dos impactos políticos, econômicos, sociais, ambientais e culturais inerentes às atividades de Engenharia.
	\item Atuação profissional comprometida com a ética, a responsabilidade social e a promoção do desenvolvimento sustentável.
\end{itemize}